%
% menu_lifecycle_bounded_inbox.tex
%
% A Z specification of the bounded per-connection Agent-Subscribe inbox and the
% leaf lock that serializes its drop-oldest overflow -- the concurrency axis
% docs/menu_lifecycle.tex and docs/menu_lifecycle_delivery.tex leave abstract.
% Bead lux-wk3p (branch chore/lux-m3xr-quality-followups); reviews the leaf-lock
% fix gvr is building for src/punt_lux/domain/hub/inbox_queue.py's BoundedInbox.
%
% WHY THIS MODEL EXISTS. lux-wk3p bounded the inbox: BoundedInbox.put drops the
% oldest undelivered message when the queue is at capacity (INBOX_CAPACITY = 32),
% a backstop against a subscriber whose recv leg has stopped draining. Review of
% that change found put's overflow logic is a non-atomic check-then-act,
%
%     if self._queue.qsize() >= self._capacity:   # (1) check
%         self._drop_oldest()                      # (2) drop the front
%     self._queue.put(message)                     # (3) enqueue
%
% run by TWO producers that do NOT share a lock:
%
%   * PUBLISH fan-out -- Hub.publish -> the session's Hub writer (inbox.py's
%     _writer) -> inbox_for(cid).put(msg). inbox_for takes the registry lock
%     _inboxes_lock only to look the inbox up, then RELEASES it; the put runs
%     with no lock held.
%   * MENU-CLICK offer -- MenuEventRouter.deliver -> inbox.offer(cid, msg), whose
%     get-then-put runs INSIDE _inboxes_lock (the D1 hold of menu_lifecycle.tex).
%
% So one producer runs put outside _inboxes_lock and the other inside it: the
% registry lock does not serialize them against each other. Their check-then-act
% puts interleave, and the compound (1)(2)(3) is torn:
%
%   OVERSHOOT. Two puts arrive at depth = cap - 1. Both read qsize < cap at (1),
%     so neither drops; both enqueue at (3). depth ends cap + 1 -- the bound the
%     inbox exists to hold is broken.
%   DOUBLE-DROP. Two puts arrive at depth = cap. Both read qsize >= cap at (1),
%     so both will drop. The first drops (depth cap -> cap-1) then the second
%     drops (cap-1 -> cap-2) -- a get_nowait fired while the queue was no longer
%     full, discarding a message that had room. Two messages discarded where one
%     over-capacity admission calls for exactly one.
%
% THE FIX gives BoundedInbox its OWN internal lock, held across the whole
% compound put (and across _drop_oldest, get, drain, depth). It is a fresh lock,
% distinct from the registry's _inboxes_lock. The implementer's earlier
% "no new lock" self-assessment was wrong: this IS a new lock, and the repo's
% z-spec policy requires deadlock-freedom to be PROVEN, not assumed, for any lock
% change -- especially one that nests.
%
% WHAT THIS MODEL PROVES, over a bounded carrier and every interleaving ProB
% reaches (two producers, capacity 2):
%
%   B1 -- bound integrity: depth never exceeds capacity.
%   B2 -- drop minimality (no double-drop): a drop fires only when the inbox is
%     genuinely full; a single over-capacity admission drops exactly one oldest.
%   DEADLOCK-FREEDOM: the two locks form a one-way acquisition edge
%     regLock -> boxLock (the offer path), the publish path takes boxLock with no
%     outer lock, and no path takes regLock while holding boxLock. The
%     acquisition graph gains only edges pointing one way; a cycle needs two. This
%     is the same one-direction-edge structure menu_lifecycle.tex proved safe for
%     the D3 StoreLock -> menu-registry-lock edge.
%
% and it reproduces both defects as fidelity controls that a model too abstract
% to trust could not:
%
%   menu_lifecycle_bounded_inbox_nolock_buggy.tex -- remove the inbox lock, so
%     the compound put's check/drop/enqueue interleave across the two producers.
%     B1's negation (depth > cap) and B2's negation (a drop at depth < cap) both
%     become reachable -- overshoot AND double-drop, the exact review finding.
%   menu_lifecycle_bounded_inbox_deadlock_buggy.tex -- add the reverse edge: the
%     publish path acquires regLock WHILE holding boxLock. The one-way property
%     is destroyed and ProB's deadlock check FINDS the cycle.
%
% Read first: src/punt_lux/domain/hub/inbox_queue.py (BoundedInbox.put,
% _drop_oldest), inbox.py (offer -- put under _inboxes_lock; _writer -- put with
% _inboxes_lock released; inbox_for, drop_session, inbox_depth_for, drain_inbox),
% and the companions docs/menu_lifecycle.tex (M1, MO) and
% docs/menu_lifecycle_delivery.tex (M1b), whose delivery invariants this change
% must not disturb (see the M1/M1b section below).
%
% Type-check:  fuzz docs/menu_lifecycle_bounded_inbox.tex
% Model-check (deadlock-freedom + the structural invariant over the state space):
%   probcli docs/menu_lifecycle_bounded_inbox.tex -model_check -p DEFAULT_SETSIZE 2
% The safety goals (B1, B2) are reachability checks; see the Verification section.
%

\documentclass[a4paper,10pt,fleqn]{article}
\usepackage[margin=1in]{geometry}
\usepackage{fuzz}
\usepackage[colorlinks=true,linkcolor=blue,citecolor=blue,urlcolor=blue]{hyperref}

\begin{document}

\title{lux Bounded Agent-Subscribe Inbox: the Leaf-Lock Refinement}
\author{Proving the drop-oldest overflow holds its bound and drops exactly once
  under two concurrent producers, and that the inbox lock is a deadlock-free leaf}
\date{September 2026}
\maketitle

% ============================================================
\section{Introduction}
% ============================================================

The Agent-Subscribe inbox is a bounded FIFO. When a full inbox admits a new
message the oldest undelivered one is dropped, so a subscriber whose
\texttt{recv} leg has stopped draining cannot grow the queue without bound. Two
producers write to that inbox, and they do not share a lock: the \emph{publish}
fan-out runs \texttt{BoundedInbox.put} with the registry lock
\texttt{\_inboxes\_lock} already released, while the \emph{menu-click}
\texttt{offer} runs its \texttt{put} with \texttt{\_inboxes\_lock} held. The
registry lock therefore serializes neither producer against the other, and
\texttt{put}'s overflow logic --- read the size, drop the front if full, enqueue
--- is a non-atomic check-then-act two producers can tear.

This model states, over a bounded carrier and every interleaving ProB reaches,
the two safety properties the bounded inbox must hold, proves the leaf lock the
fix adds restores them, and proves that lock introduces no deadlock.

\begin{description}
  \item[B1 --- bound integrity.] The inbox depth never exceeds its capacity:
    \[ depth \leq cap. \]
    This is what the un-serialized check-then-act breaks by \emph{overshoot}:
    two puts at $depth = cap - 1$ both read ``not full'' and both enqueue,
    leaving $depth = cap + 1$.
  \item[B2 --- drop minimality (no double-drop).] A message is dropped only from
    a genuinely full inbox; a drop never discards a message that had room:
    \[ \lnot\, (spuriousDrop = fset). \]
    This is what the un-serialized check-then-act breaks by \emph{double-drop}:
    two puts at $depth = cap$ both decide to drop, and the second
    \texttt{get\_nowait} fires after the first has already made room --- a drop
    at $depth < cap$, one over-capacity admission costing two discarded messages.
\end{description}

Alongside the two safety properties, the model carries the two locks explicitly
and proves the third obligation the repo's z-spec policy attaches to any lock
change:

\begin{description}
  \item[Deadlock-freedom.] The inbox lock (\texttt{boxLock}) is a leaf. The only
    nesting is \texttt{regLock} $\to$ \texttt{boxLock}, on the offer path; the
    publish path takes \texttt{boxLock} with no outer lock; no path takes
    \texttt{regLock} while holding \texttt{boxLock}. The acquisition graph gains
    only edges pointing one way, and a cycle needs two. A full model-check finds
    no deadlock.
\end{description}

As in the companions, the safety properties are kept out of the state schema:
a predicate placed in the state invariant is enforced as a guard on every
successor, silently disabling the very operation that would break it, so a
violating operation becomes unreachable rather than a counterexample. The schema
below carries only structural coherence; B1 and B2 are checked as the
reachability of their negations (Section~\ref{sec:verify}), and deadlock-freedom
by the full model-check.

% ============================================================
\section{Constants}
% ============================================================

The capacity is bounded small; two suffices to exhibit both the overshoot at
$cap - 1$ and the double-drop at $cap$. \texttt{maxDepth} is one above capacity,
the least bound that can \emph{represent} an overshoot --- the fixed model never
reaches it, and the no-lock control does.

\begin{axdef}
  cap : \nat \\
  maxDepth : \nat
\where
  cap = 2 \\
  maxDepth = 3
\end{axdef}

% ============================================================
\section{Free Types}
% ============================================================

\begin{zed}
  FSTATE ::= fclear | fset
\end{zed}

A two-value flag, matching the companions' \texttt{FSTATE}. Here \texttt{spuriousDrop}
carries it: whether any \texttt{\_drop\_oldest} has fired while the inbox was not
full.

\begin{zed}
  REGH ::= rfree | roffer
\end{zed}

Who holds the registry lock \texttt{\_inboxes\_lock}. In the fixed model only the
offer path holds it (the publish path's transient hold inside \texttt{inbox\_for}
is released before \texttt{put} runs, so it never coexists with \texttt{boxLock}
and is not modelled); the reverse-edge control adds a publish holder.

\begin{zed}
  BOXH ::= bfree | boffer | bpublish
\end{zed}

Who holds the inbox's own lock \texttt{boxLock} --- free, the offer producer, or
the publish producer. At most one holder: this is the mutual exclusion that makes
the compound \texttt{put} atomic.

\begin{zed}
  OPHASE ::= oidle | oreg | oboth
\end{zed}

Where the offer producer is in its two-lock climb: idle, holding
\texttt{regLock} and waiting for \texttt{boxLock} (\texttt{oreg} --- the blockable
wait that a reverse edge would turn into half a deadlock), or holding both and
ready to put.

\begin{zed}
  PPHASE ::= pidle | pbox
\end{zed}

Where the publish producer is: idle, or holding \texttt{boxLock} and ready to
put. It never waits for a second lock --- the property the deadlock proof turns
on.

% ============================================================
\section{State}
% ============================================================

One flat schema. \texttt{depth} is the queued-but-undelivered count
(\texttt{qsize}); \texttt{spuriousDrop} records whether any drop fired below
capacity. \texttt{regLock} and \texttt{boxLock} are the two lock holders, and the
two phases place each producer in its lock climb. The at-most-one-holder
discipline is structural: a lock is held by exactly one value of its holder type.

\begin{schema}{InboxLock}
  depth : \nat \\
  spuriousDrop : FSTATE \\
  regLock : REGH \\
  boxLock : BOXH \\
  oPhase : OPHASE \\
  pPhase : PPHASE
\where
  depth \leq maxDepth \\
  (oPhase = oidle \iff regLock = rfree) \\
  (oPhase = oboth \implies boxLock = boffer) \\
  (pPhase = pbox \iff boxLock = bpublish)
\end{schema}

The coherence predicates tie the phases to the locks: the offer producer holds
\texttt{regLock} exactly when it is past \texttt{oidle}, holds \texttt{boxLock}
when it reaches \texttt{oboth}, and the publish producer holds \texttt{boxLock}
exactly when it is in \texttt{pbox}. \texttt{depth} $\leq$ \texttt{maxDepth} is a
carrier bound, not the safety property: it says the count fits the modelled
range, and it is deliberately one \emph{above} \texttt{cap} so an overshoot is
representable. B1, the real bound, is checked as reachability.

% ============================================================
\section{Initialization}
% ============================================================

\begin{schema}{Init}
  InboxLock'
\where
  depth' = 0 \\
  spuriousDrop' = fclear \\
  regLock' = rfree \\
  boxLock' = bfree \\
  oPhase' = oidle \\
  pPhase' = pidle
\end{schema}

An empty inbox, both locks free, both producers idle, nothing dropped.

% ============================================================
\section{Operations}
% ============================================================

The compound \texttt{put} appears as one atomic operation per producer
(\texttt{OfferPut}, \texttt{PubPut}), each performed while that producer holds
\texttt{boxLock}. Representing a lock-protected critical section as a single Z
operation is the standard abstraction, and it is exactly how
\texttt{menu\_lifecycle.tex} models the D1 fix: because \texttt{boxLock} is
mutually exclusive and held across the whole compound, no other producer's put
can interleave the check, the drop, and the enqueue. The no-lock control drops
that hold and splits the compound into its three tearable steps.

\subsection{OfferAcqReg --- the offer path takes the registry lock}

Models entering \texttt{offer}'s \texttt{with \_inboxes\_lock} block. The offer
producer takes \texttt{regLock} and advances to \texttt{oreg}, where it will wait
for \texttt{boxLock}.

\begin{schema}{OfferAcqReg}
  \Delta InboxLock \\
\where
  oPhase = oidle \\
  regLock = rfree \\
  regLock' = roffer \\
  oPhase' = oreg \\
  depth' = depth \\
  spuriousDrop' = spuriousDrop \\
  boxLock' = boxLock \\
  pPhase' = pPhase
\end{schema}

\subsection{OfferAcqBox --- the offer path takes the inbox lock, nested}

Models \texttt{inbox.put} entering the inbox's own lock \emph{while
\texttt{\_inboxes\_lock} is still held}. This is the one nesting in the design:
\texttt{regLock} $\to$ \texttt{boxLock}. The step is guarded on
\texttt{boxLock} free, so if the publish producer holds it the offer producer
waits here in \texttt{oreg} --- the blockable wait.

\begin{schema}{OfferAcqBox}
  \Delta InboxLock \\
\where
  oPhase = oreg \\
  boxLock = bfree \\
  boxLock' = boffer \\
  oPhase' = oboth \\
  depth' = depth \\
  spuriousDrop' = spuriousDrop \\
  regLock' = regLock \\
  pPhase' = pPhase
\end{schema}

\subsection{OfferPut --- the compound put, atomic under boxLock}

Models the whole of \texttt{BoundedInbox.put} running under the inbox lock: if
the inbox is full one message is dropped and one enqueued (depth unchanged); else
one is enqueued (depth up one). The drop, when it fires, fires at $depth = cap$,
so it never lowers a queue that had room --- \texttt{spuriousDrop} stays clear.
Both locks release as the two \texttt{with} blocks close.

\begin{schema}{OfferPut}
  \Delta InboxLock \\
\where
  oPhase = oboth \\
  (depth < cap \implies depth' = depth + 1) \\
  (depth \geq cap \implies depth' = depth) \\
  spuriousDrop' = spuriousDrop \\
  boxLock' = bfree \\
  regLock' = rfree \\
  oPhase' = oidle \\
  pPhase' = pPhase
\end{schema}

\subsection{PubAcqBox --- the publish path takes the inbox lock, un-nested}

Models the publish \texttt{\_writer}'s \texttt{inbox\_for(cid).put}. The registry
lock \texttt{inbox\_for} took to look the inbox up has already been released, so
the publish producer takes \texttt{boxLock} with \emph{no} outer lock held. It
never climbs to a second lock --- there is no \texttt{PubAcqReg} in the fixed
model.

\begin{schema}{PubAcqBox}
  \Delta InboxLock \\
\where
  pPhase = pidle \\
  boxLock = bfree \\
  boxLock' = bpublish \\
  pPhase' = pbox \\
  depth' = depth \\
  spuriousDrop' = spuriousDrop \\
  regLock' = regLock \\
  oPhase' = oPhase
\end{schema}

\subsection{PubPut --- the publish compound put, atomic under boxLock}

The same atomic compound as \texttt{OfferPut}, under the publish producer's hold
of \texttt{boxLock}. Releases only \texttt{boxLock}: the publish path never held
\texttt{regLock}.

\begin{schema}{PubPut}
  \Delta InboxLock \\
\where
  pPhase = pbox \\
  (depth < cap \implies depth' = depth + 1) \\
  (depth \geq cap \implies depth' = depth) \\
  spuriousDrop' = spuriousDrop \\
  boxLock' = bfree \\
  pPhase' = pidle \\
  regLock' = regLock \\
  oPhase' = oPhase
\end{schema}

\subsection{Drain --- recv consumes one message}

Models a subscriber's \texttt{recv} (\texttt{get}) draining one message. In the
fix \texttt{get} also runs under the inbox lock, so \texttt{Drain} is guarded on
\texttt{boxLock} free --- it cannot run while a producer holds the lock, and
holds no lock in a wait state of its own, so it contributes no deadlock edge. It
lets \texttt{depth} fall so contention recurs and B1/B2 are not vacuous.

\begin{schema}{Drain}
  \Delta InboxLock \\
\where
  boxLock = bfree \\
  depth > 0 \\
  depth' = depth - 1 \\
  spuriousDrop' = spuriousDrop \\
  regLock' = regLock \\
  boxLock' = boxLock \\
  oPhase' = oPhase \\
  pPhase' = pPhase
\end{schema}

% ============================================================
\section{Invariants}
\label{sec:inv}
% ============================================================

\paragraph{B1 --- bound integrity.}
\[ depth \leq cap. \]
No operation can push \texttt{depth} past \texttt{cap}. \texttt{OfferPut} and
\texttt{PubPut} each require their producer to hold \texttt{boxLock}, and
\texttt{boxLock} admits one holder, so the two compounds never interleave: each
runs atomically against a \texttt{depth} it alone is changing. From
$depth \leq cap$ an atomic compound leaves $depth + 1 \leq cap$ (if it was below)
or $depth$ (if it was full), so $depth \leq cap$ is preserved. \texttt{maxDepth}
is reachable only if the bound breaks; against the fixed model $depth = maxDepth$
is not found. The no-lock control removes the \texttt{boxLock} hold, and
$depth = maxDepth$ becomes reachable at once.

\paragraph{B2 --- drop minimality.}
\[ \lnot\, (spuriousDrop = fset). \]
A drop fires only inside an atomic compound, and only on the $depth \geq cap$
branch, so it only ever discards the oldest message of a genuinely full inbox.
\texttt{spuriousDrop} is never set in the fixed model. The no-lock control lets
one producer's drop-decision go stale --- taken at $depth = cap$, executed after
another producer has already dropped, at $depth < cap$ --- and
\texttt{spuriousDrop} becomes reachable: two discards for one over-capacity
admission.

\paragraph{Deadlock-freedom and the one-way lock edge.}
\label{sec:lockorder}
The inbox lock is new, and it nests: \texttt{offer} holds \texttt{regLock}
(\texttt{\_inboxes\_lock}) across \texttt{inbox.put}, which now takes
\texttt{boxLock} --- the edge \texttt{regLock} $\to$ \texttt{boxLock}. Every other
holder of \texttt{boxLock} takes it alone: the publish \texttt{\_writer} runs
\texttt{put} after \texttt{inbox\_for} has released \texttt{regLock};
\texttt{inbox\_depth\_for} and \texttt{drain\_inbox} read the inbox out under
\texttt{regLock}, release it, and only then take \texttt{boxLock} for
\texttt{depth}/\texttt{drain}; and \texttt{drop\_session} takes \texttt{regLock}
alone and never touches \texttt{boxLock} at all. No path takes \texttt{regLock}
while holding \texttt{boxLock}. The acquisition graph is the single edge
\texttt{regLock} $\to$ \texttt{boxLock}; there is no reverse edge, so no cycle,
so no deadlock. This is the same shape \texttt{menu\_lifecycle.tex}'s
Section~``lock order'' establishes for the D3 edge \texttt{StoreLock} $\to$
menu-registry lock: one more caller taking two locks in the one order every
caller already uses. The model carries the two locks and the offer producer's
blockable \texttt{oreg} wait so a full model-check exercises the contention; it
finds no deadlock. The reverse-edge control adds the missing back edge --- the
publish path acquiring \texttt{regLock} while holding \texttt{boxLock} --- and the
deadlock check FINDS the cycle, so the one-way property is shown load-bearing,
not assumed.

% ============================================================
\section{Fidelity Controls}
\label{sec:fidelity}
% ============================================================

Each control differs from this specification in one place and makes an
obligation's negation reachable where the fixed spec makes it unreachable.

\paragraph{No lock --- \texttt{menu\_lifecycle\_bounded\_inbox\_nolock\_buggy.tex}.}
\label{sec:fidelity-nolock}
The inbox lock is removed, so the compound \texttt{put} is no longer atomic: it
splits into \texttt{PutCheck} (read \texttt{depth}, decide whether to drop),
\texttt{PutDrop} (the \texttt{get\_nowait}), and \texttt{PutEnq} (the enqueue),
one triple per producer, freely interleaved. Two traces make the negations
reachable. \emph{Overshoot:} from $depth = cap - 1$, both producers
\texttt{PutCheck} (neither will drop), then both \texttt{PutEnq} ---
$depth = cap + 1 = maxDepth$, B1's negation, FOUND. \emph{Double-drop:} from
$depth = cap$, both \texttt{PutCheck} (both will drop), the first \texttt{PutDrop}
lowers $depth$ to $cap - 1$, the second \texttt{PutDrop} then fires at
$depth < cap$ --- \texttt{spuriousDrop} $=$ \texttt{fset}, B2's negation, FOUND.
The fix (atomic compound under \texttt{boxLock}) removes both windows: no second
producer can interleave a check, a drop, and an enqueue that are one step.

\paragraph{Reverse edge --- \texttt{menu\_lifecycle\_bounded\_inbox\_deadlock\_buggy.tex}.}
\label{sec:fidelity-deadlock}
The publish path is made to acquire \texttt{regLock} \emph{while holding}
\texttt{boxLock} --- a hypothetical in which \texttt{\_writer} reached back into
the registry lock under the inbox lock. The one nesting \texttt{regLock} $\to$
\texttt{boxLock} now has its reverse, \texttt{boxLock} $\to$ \texttt{regLock}, and
the acquisition graph has a cycle. The trace
\texttt{OfferAcqReg}\,;\,\texttt{PubAcqBox} reaches the state where the offer
producer holds \texttt{regLock} and waits for \texttt{boxLock} while the publish
producer holds \texttt{boxLock} and waits for \texttt{regLock}: no operation is
enabled --- a deadlock, FOUND by ProB. This is the counterfactual the fixed
model's one-way edge rules out, shown load-bearing.

% ============================================================
\section{The Delivery Invariants M1 and M1b Are Preserved}
\label{sec:m1}
% ============================================================

The bounded inbox and its lock do not disturb the delivery properties the
companions prove. Three facts carry that.

First, \textbf{M1} (\texttt{menu\_lifecycle.tex}, no false-positive delivery) is
the orphan-put property: \texttt{offer} must not report \texttt{delivered} for a
put into an inbox \texttt{drop\_session} has already removed. Its guarantee is
\texttt{offer} holding \texttt{\_inboxes\_lock} (\texttt{regLock}) across its
\texttt{get} and \texttt{put}, serialized against \texttt{drop\_session}'s hold of
the same lock. The new lock is \texttt{boxLock}, \emph{internal to} the
\texttt{BoundedInbox} the \texttt{get} returned; it neither loosens nor moves the
\texttt{regLock} hold M1 rests on. M1 is unaffected.

Second, \textbf{M1b} (\texttt{menu\_lifecycle\_delivery.tex}, no live-recipient
loss) is the property that no \emph{undrained} event is discarded from a
still-registered connection by the departure race --- guaranteed by the departure
removing the connection from the registry before the cascade tears the inbox down.
The bounded inbox changes neither the departure ordering nor the registry-discard
step, and the second producer (publish) only ever calls \texttt{put}: it moves no
departure and discards nothing on teardown. M1b's race-freedom is unaffected.

Third, the bounded drop is a \emph{separate, intentional} loss policy, not an M1b
violation. \texttt{\_drop\_oldest} discards the oldest undelivered message of a
full inbox belonging to a possibly still-registered but non-draining subscriber
--- by design, the overflow backstop \texttt{INBOX\_CAPACITY} exists to apply, the
same policy the callback hold's \texttt{\_HOLD\_CAPACITY} applies. M1b forbids
loss from the delivery/departure \emph{race}; it does not forbid the capacity
backstop, which is the documented behaviour of a bounded queue and is what B1/B2
above govern. Reading the capacity drop as an M1b breach would be a category
error: M1b is about races, the bound is about overflow.

Because the change touches none of the transitions \texttt{menu\_lifecycle.tex}
and \texttt{menu\_lifecycle\_delivery.tex} model, those specs continue to hold
unchanged; re-running their \texttt{fuzz} and goal checks confirms M1 and M1b
still pass with the bound and the lock in place.

% ============================================================
\section{Verification}
\label{sec:verify}
% ============================================================

Type-checking is a \texttt{fuzz} pass:
\begin{verbatim}
  fuzz docs/menu_lifecycle_bounded_inbox.tex
\end{verbatim}

B1 and B2 are safety properties, checked by asking ProB whether their negation
is reachable. A goal reported \emph{not found} means the invariant holds on every
reachable state.

\begin{verbatim}
  # B1 -- bound integrity (goal NOT found in this, the fixed, spec)
  probcli docs/menu_lifecycle_bounded_inbox.tex -model_check -nodead \
    -goal "depth > cap" \
    -p DEFAULT_SETSIZE 2

  # B2 -- drop minimality (goal NOT found in this, the fixed, spec)
  probcli docs/menu_lifecycle_bounded_inbox.tex -model_check -nodead \
    -goal "spuriousDrop = fset" \
    -p DEFAULT_SETSIZE 2
\end{verbatim}

The positive outcomes are confirmed reachable, so the invariants are not
vacuously true of a model that can never fill or drain the inbox:

\begin{verbatim}
  # the inbox reaches capacity (goal FOUND)
  probcli docs/menu_lifecycle_bounded_inbox.tex -model_check -nodead \
    -goal "depth = cap" \
    -p DEFAULT_SETSIZE 2

  # both producers can hold the lock over the run, one at a time (goal FOUND)
  probcli docs/menu_lifecycle_bounded_inbox.tex -model_check -nodead \
    -goal "boxLock = bpublish" \
    -p DEFAULT_SETSIZE 2
\end{verbatim}

Deadlock-freedom and the structural state invariant are checked by a full
model-check over the reachable state space (deadlock detection is on by default):

\begin{verbatim}
  probcli docs/menu_lifecycle_bounded_inbox.tex -model_check -p DEFAULT_SETSIZE 2
\end{verbatim}

The fidelity controls make the negations \emph{FOUND}:

\begin{verbatim}
  # No lock: B1's negation FOUND (overshoot)
  probcli docs/menu_lifecycle_bounded_inbox_nolock_buggy.tex -model_check -nodead \
    -goal "depth > cap" \
    -p DEFAULT_SETSIZE 2

  # No lock: B2's negation FOUND (double-drop)
  probcli docs/menu_lifecycle_bounded_inbox_nolock_buggy.tex -model_check -nodead \
    -goal "spuriousDrop = fset" \
    -p DEFAULT_SETSIZE 2

  # Reverse edge: the deadlock check FINDS a deadlock
  probcli docs/menu_lifecycle_bounded_inbox_deadlock_buggy.tex -model_check \
    -p DEFAULT_SETSIZE 2
\end{verbatim}

Re-run \texttt{fuzz} and these goal checks whenever \texttt{inbox\_queue.py}'s
\texttt{put}/\texttt{\_drop\_oldest} or its lock, or \texttt{inbox.py}'s
\texttt{offer}/\texttt{\_writer}/\texttt{inbox\_for}/\texttt{drop\_session}
(the lock nesting), change. Re-run the companions' checks too: M1 and M1b are
proved there, and this change is designed to leave them untouched.

\end{document}
